\section{명제를 읽기 위한 형식 언어(EasyCrypt)}
\label{sec:formal-language}

\subsection{핵심 술어와 관계}

\begin{definition}[패킹된 키(packed-key) 접두부]
\label{def:vk-prefix}
\[
  \operatorname{VKPrefix}_{n}(sk,vk)
  \;\Longleftrightarrow\;
  \forall i\;(0\le i<n\Rightarrow sk[i]=vk[i]).
\]
코드의 \identifier{vk_prefix_eq}와 같으며, mode~2에서는 \(n=992\)이다.
\end{definition}

\begin{definition}[외부 키(key) 영역의 접두부 일치]
\label{def:external-prefix}
\[
  \operatorname{ExternalKeyPrefix}(mem,sku,vku)
  \Longleftrightarrow
  \forall i\in[0,992),\;
  \operatorname{load8}(mem,sku+i)=\operatorname{load8}(mem,vku+i).
\]
코드의 \identifier{external_key_prefix_match}이며, 실제 raw KeyGen이 서로 다른
외부 SK/VK 주소에 내보낸 첫 992바이트를 비교한다.
\end{definition}

\begin{definition}[메모리에서의 접두부]
\label{def:mem-prefix}
\[
  \operatorname{SKMemPrefix}(sk,mem,b)
  \;\Longleftrightarrow\;
  \forall i\in[0,992),\;
  sk[i]=\operatorname{load8}(mem,\operatorname{uint}(b)+i).
\]
이 술어는 Sign의 로컬 \identifier{BArray2752}와 Verify의 raw-memory 읽기를
한 관계식 안에서 비교한다.
\end{definition}

\begin{definition}[정준 주소 구간(canonical address range)]
\label{def:canonical-region}
정수 주소 \(a\)와 길이 \(n\)에 대해
\[
\begin{aligned}
  \operatorname{CanonicalAddress}(a)
    &\Longleftrightarrow 0\le a<2^{64},\\
  \operatorname{CanonicalRegion}(a,n)
    &\Longleftrightarrow
      0\le a\;\land\;0\le n\;\land\;a+n\le 2^{64}.
\end{aligned}
\]
첫 전제는 원시 응용 이진 인터페이스(raw application binary interface, raw
ABI)의 정수 주소가 \identifier{W64.of_int}와
\identifier{W64.to_uint}를 왕복하게 하고, 둘째는 주소 구간의 모듈러
순환(wraparound)을 배제한다. 길이가 0이면 \(a=2^{64}\)도 둘째 식만으로는
허용되므로, 실제
왕복 정리는 필요한 주소와 길이에 \identifier{canonical_ui64_address}를
별도로 요구한다. 이것은 물리적 메모리 안전성 증명이 아니라 호출 계약이다.
\end{definition}

\begin{definition}[메모리 구간 등식]
\label{def:region-eq}
\[
  \operatorname{RegionEq}(m_1,m_2,b,n)
  \Longleftrightarrow
  \forall i\in[0,\operatorname{uint}(n)),\;
  \operatorname{load8}(m_1,\operatorname{uint}(b)+i)
  =\operatorname{load8}(m_2,\operatorname{uint}(b)+i).
\]
코드의 \identifier{region_eq}다. Week~6의 hash 정리는 전체 memory 등식 대신
VK, pre, message에서 실제로 읽는 구간의 이 관계만 요구한다.
\end{definition}

\begin{definition}[정준 서명 접두부(canonical signature prefix) 입력]
\label{def:signature-prefix-canonical}
\identifier{canonical_challenge}는 256개 challenge word가 각각 정확히 0 또는
1임을 뜻한다. \identifier{canonical_signed_low}는 1024개 low word 각각이
자신의 하위 바이트를 부호확장한 값이고 signed 정수로 \([-128,128)\)에 있음을
뜻한다. 실제 prefix round trip은 이 두 전제 아래 성립한다.
\end{definition}

\begin{definition}[정준 mode-2 HBZ와 잎 프레임(leaf frame)]
\label{def:canonical-hbz}
용량 \(2048\)의 32비트 계수 배열 \(h\)에 대해
\[
  \operatorname{CanonicalHBZ}_2(h)
  \Longleftrightarrow
  \forall i\in[0,1024),\;-6\le\operatorname{sint}(h[i])<7.
\]
코드의 \identifier{canonical_hbz_mode2}다.
\identifier{prepared_hbz_prefix}는 첫 \(n\)개 symbol byte가
\(\operatorname{sint}(h[i])+6\)임을, \identifier{decoded_hbz_prefix}는 첫
\(n\)개 복원 계수가 원본과 같음을 뜻한다. \identifier{byte_tail_frame}과
\identifier{coeff_tail_frame}은 각각 \([1024,2048)\)이 leaf 실행 전후에
변하지 않음을 기록한다.
\end{definition}

\begin{definition}[mode-2 HBZ 표 인증서(table certificate)]
\label{def:hbz-table-certificate}
\identifier{mode2_hbz_table_certificate}는 다음 세 사실의 conjunction이다.
\begin{enumerate}
  \item 13개 symbol의 실제 encoder table 네 field가 명세식과 같다.
  \item 실제 decoder word가 각 interval의 start/frequency를 정확히 담는다.
  \item 실제 512개의 packed symbol word가 1024개 slot 각각을 유일한
        interval symbol에 연결한다.
\end{enumerate}
\identifier{actual_mode2_hbz_tables_certified}는 pack/unpack 추출물의 literal
배열에 이 술어를 적용한 정리다. certificate를 생성한 외부 스크립트 자체를
가정하는 것이 아니라, 생성된 모든 EasyCrypt 등식을 kernel이 다시 검사한다.
\end{definition}

\begin{definition}[정준 rANS 바이트 추적(canonical rANS byte trace)]
\label{def:rans-byte-trace}
정준 기호 목록(canonical symbol list) \(s=(s_0,\ldots,s_{1023})\)에 대해
\identifier{encode_trace}는 \(x_{1024}=2^{23}\)에서 기호(symbol)를 역순으로
처리해 경계 상태(boundary state) \((x_j)_{j=0}^{1024}\)와 디코더 순서 바이트
구간(decoder-order byte segment) \((\beta_j)_{j=0}^{1023}\)를 결정한다. 커서
절단값(cursor cut)과 전체 바이트열(byte sequence)은
\[
  c_0=4,\qquad c_{j+1}=c_j+|\beta_j|,
  \qquad
  B=\operatorname{serialize32\_le}(x_0)
      \Vert\beta_0\Vert\cdots\Vert\beta_{1023}
\]
이다. 코드의 \identifier{valid_rans_trace}(s,xs,cuts,bytes)는 단지
\identifier{trace_states}, \identifier{trace_cuts}, \identifier{trace_bytes}로
계산한 이 세 객체와 인자가 같다는 논리곱(conjunction)이다. 외부 추적(external
trace)의 존재나 디코더 성공(decoder success)을 가정하는 술어가 아니며,
\identifier{valid_rans_trace_canonical}이 정준 구성(canonical construction)을
증명한다.
\end{definition}

\begin{definition}[배열/목록 연결과 구간 프레임(array/list bridge and segment frame)]
\label{def:rans-array-list-bridge}
용량 2048의 바이트 배열 \(a\)에 대해
\identifier{symbol_list_of_array}(a)는 앞 1024개 원소의 정수 목록이고,
\identifier{symbol_suffix}(a,i)는 그 목록에서 앞 \(i\)개를 버린 접미부다.
\identifier{segment_matches}(a,p,bs)는
\[
  0\le p\;\land\;p+|bs|\le2048\;\land\;
  \forall k\in[0,|bs|),\quad
  \operatorname{uint}(a[p+k])=bs[k]
\]
를 뜻한다. \identifier{prefix_frame}(before,after,p)는 모든 \(k<p\)에서 두
배열의 \(k\)번째 바이트가 같다는 뜻이다. Week~10의 컴파일된 정리는
\[
\begin{aligned}
  \mathsf{symbol\_suffix}(a,1024)&=[],\\
  \mathsf{symbol\_suffix}(a,i)
    &=\operatorname{uint}(a[i])::\mathsf{symbol\_suffix}(a,i+1)
\end{aligned}
\]
와 구간의 이어 붙이기, 앞쪽 한 바이트 쓰기, 프레임 축소 규칙을 제공한다.
\end{definition}

\begin{definition}[실제 인코더 내부 단계 불변식(actual encoder inner-step invariant)]
\label{def:rans-actual-inner}
\identifier{encoder_inner_phase_inv}와
\identifier{encoder_inner_segment_inv}는 생성된
\identifier{_rans_encode}의 내부 정규화 반복문을 직접 연다. 진입 상태
\((x_0,off_0)\)와 단계 \(k\)가 존재하여
\[
\begin{gathered}
  4\le off_0\le1024,\qquad 0\le k\le4,
  \qquad 1\le\operatorname{uint}(x_{\max}),\\
  \operatorname{uint}(off)=off_0-k,\qquad
  \operatorname{uint}(x)=\mathsf{encoder\_phase\_state}(x_0,k),\\
  \mathsf{segment\_matches}
    (encp,off_0-k,\mathsf{encoder\_phase\_written}(x_0,k))
\end{gathered}
\]
이고 해당 내부 호출의 진입 배열보다 낮은 접두부가 보존된다는 뜻이다. 이는
Week~10의 보수적 국소 술어로서 실제 한 바이트 저장과 W64 커서 안전성을
검증한다. Week~11의 더 강한 꼬리 인식 술어가 외부 누적과 정확한 0/1/2바이트
길이를 추가한다.
\end{definition}

\begin{definition}[꼬리 인식 내부 불변식(tail-aware inner invariant)]
\label{def:rans-tail-aware-inner}
Week~11의 \identifier{encoder_inner_tail_inv}는 정확한 순수 정규화 길이
\(\ell=\mathsf{mode2\_normalization\_len}(x_0,s)\)와 이미 존재하는 순수
기호 꼬리 \(tail\)의 인코딩 추적을 함께 둔다. 핵심 구간 조건은
\[
  \mathsf{segment\_matches}\bigl(
    encp,off_0-k,
    \mathsf{inner\_written\_suffix}(x_0,s,k)\mathbin{++}
      \mathsf{encode\_trace}(tail).\mathsf{bytes}
  \bigr)
\]
이고 \(0\le k\le\ell\), 커서 등식, 상태 등식과 접두부 프레임을 함께
요구한다. 실제 가드(actual guard)의 진행과 종료에서 이 술어가 보존되고 정확한
\(k=\ell\le2\)가 도출된다. 이어 외부 한 기호 점화식과 최종 저장 합성에 사용되어
실제 외부 절차의 전체 성공 사후조건을 닫는다.
\end{definition}

\begin{verified}[실제 인코더의 성공 조건부 전체 접미부]
\label{thm:rans-encoder-suffix}
\identifier{actual_rans_encode_trace_closure}는 실제 생성
\identifier{_rans_encode}를 직접 열고, 반환 \(bad\ne0\)이거나 \(bad=0\)일 때
\[
  \mathsf{segment\_matches}
  (encp',\operatorname{uint}(off'),
   \mathsf{trace\_bytes}(\mathsf{symbol\_list\_of\_array}(symbols_0)))
\]
와 \(off'+|\mathsf{trace\_bytes}|=1024\), 성공 크기 범위 및
\identifier{prefix_frame}이 모두 성립함을 증명한다. 성공은 사전조건이 아니다.
그러나 Hoare 부분정확성이므로 종료와 성공 도달은 별도 의무다.
\end{verified}

\begin{definition}[실제 디코더 정확 추적 입력(actual decoder exact-trace input)]
\label{def:rans-decoder-exact-input}
\identifier{actual_mode2_decoder_trace_input}(a,b,st,n)은 예상 기호 배열 \(a\),
인코딩 버퍼 \(b\), 초기 상태 \(st\)와 논리적 크기 \(n\) 사이의 함수적 표현
계약이다. 핵심은
\[
\begin{gathered}
  \mathsf{mode2\_hbz\_symbol\_stream}(a),\qquad
  n=|\mathsf{trace\_bytes}(\mathsf{symbol\_list\_of\_array}(a))|,\\
  4\le n\le1024,\qquad
  \mathsf{segment\_matches}(b,0,\mathsf{trace\_bytes}(\cdots)),\\
  (st[0],st[1],st[2])=(1024,n,13)
\end{gathered}
\]
와 모든 정규화 구간에 대한 점별 읽기 등식이다. 실제 mode-2
\identifier{symbol_words}와 \identifier{dsyms_words} 배열은 최상위 Hoare
정리의 별도 바인딩으로 고정된다. 이 술어는 \(bad=0\), 출력 배열 등식 또는
사후상태를 전제로 넣지 않는다.
\end{definition}

\begin{definition}[디코더 외부/내부 추적 불변식(decoder outer/inner trace invariant)]
\label{def:rans-decoder-trace-invariants}
\identifier{decoder_outer_trace_live}는 반복 번호 \(i\)에서 실제 상태와 커서를
\[
  x=\mathsf{decoder\_state\_at}(s,i),\qquad
  \operatorname{uint}(off)=\mathsf{decoder\_cursor}(s,i)
\]
에 묶고, \(decoded[0..i)\)의 예상 기호 일치와
\(decoded[i..2048)\)의 입력 배열 프레임을 함께 유지한다.
\identifier{decoder_inner_trace_live}는 한 기호의 단어 갱신 뒤
\[
\begin{aligned}
  k&=\operatorname{uint}(off)-\mathsf{decoder\_cursor}(s,i),\\
  x&=\mathsf{decoder\_replay\_prefix}
       (\mathsf{decoder\_state\_at}(s,i+1),s_i,k)
\end{aligned}
\]
를 두고 \(0\le k\le|\mathsf{decoder\_segment\_at}(s,i)|\)를 유지한다. 생성
가드는 정확히 \(k<|\mathsf{segment}|\)와 일치하며 구간 길이는 0, 1 또는
2바이트다. 따라서 각 실제 읽기는 한 바이트 소비 전이이거나 다음 외부 경계로
가는 무소비(no-consume) 종료다.
\end{definition}

\begin{verified}[실제 디코더의 정확 추적 복원]
\label{thm:rans-decoder-trace-refinement}
\identifier{actual_rans_decode_trace_refinement}는
\cref{def:rans-decoder-exact-input}과 실제 두 표의 바인딩 아래 실제 생성
\identifier{RansDecodeTarget.M._rans_decode}를 직접 연다. 모든 종료 실행에서
\[
  bad'=0,\qquad off'=n,\qquad
  decoded'[0..1024)=a[0..1024)
\]
이고 \([1024,2048)\) 출력 꼬리가 초기 배열과 같다. 외부 반복 출구의 내부
\(x=2^{23}\)는 마지막 실제 거절 검사를 배제하지만 ABI 사후조건에는 나타나지
않는다. 이는 정확 추적 Hoare 부분정확성이며 종료나 실제 인코더 성공 증명은
아니다.
\end{verified}

\begin{definition}[실제 rANS 결합 제어(actual rANS harness control)]
\label{def:rans-actual-harness}
\identifier{Mode2RansActualHarness.run}은 실제
\identifier{_rans_encode}를 호출한 뒤 반환된 \(bad\)가 0일 때만 실제
\identifier{__copy_encoded_suffix}와 \identifier{_rans_decode}를 호출한다.
디코더 상태의 입력 필드는 그 성공 분기에서 정확히
\((\mathsf{count},\mathsf{size},m)=(1024,\mathsf{encoded\_size},13)\)이다.
현재 결합 제어 정리는
\[
  bad_{\mathrm{enc}}\in\{0,1\},\qquad
  \mathsf{decoder\_ran}\Longleftrightarrow bad_{\mathrm{enc}}=0
\]
과 이 필드 초기화만 결론낸다. 이 제어 정리 자체는 역함수가 아니다. Week~13의
별도 의미 정리는 \cref{thm:rans-encoder-suffix}와
\cref{thm:rans-decoder-trace-refinement}, 실제 복사 정리를 이 모듈 안에서
순차 호출해 다음 결론을 닫는다.
\end{definition}

\begin{verified}[실제 인코더--복사--디코더 핵심 역함수]
\label{thm:rans-core-actual-inverse}
\identifier{actual_rans_encode_copy_decode_inverse}는 초기 인자 바인딩과
\identifier{mode2_hbz_symbol_stream}만을 전제로 한다. 종료한 실행에서 실제
인코더가 실패하면 디코더는 실행되지 않고 초기 디코더 배열/상태가 보존된다.
실제 인코더가 성공하면 복사된 버퍼가 정확한 \identifier{trace_bytes} 입력이
되고, 디코더는
\[
  bad_{\mathrm{dec}}=0,\qquad
  off_{\mathrm{dec}}=\mathsf{encoded\_size},\qquad
  decoded[0..1024)=symbols[0..1024)
\]
를 만족하며 \([1024,2048)\) 꼬리를 보존한다. 인코더 성공, 복사 구간 등식,
디코더 성공이나 출력 등식은 사전조건이 아니다. 이 정리는 Hoare
부분정확성이므로 성공 도달과 두 반복문의 종료는 별도 의무다.
\end{verified}

\begin{boundary}[성공 조건부(success-conditioned) HBZ]
\label{bdry:hbz-success}
\(\operatorname{CanonicalHBZ}_2(h)\)는 prepare leaf가 성공함을 보장하지만 고정
1024바이트 rANS workspace의 actual full encoder가 nonzero size를 반환함을
보장하지 않는다. 따라서 full round trip에서 encoder success나 decoder success를
독립된 사전조건으로 가정하면 비공허성이 사라진다. 실제 return의
\(\mathsf{size}=0\) branch를 결론에 남기거나, actual successful witness를 별도로
증명해야 한다.
\end{boundary}

\begin{definition}[\(\mu\) 접두부]
\label{def:mu-prefix}
Sign이 만드는 64바이트 배열을 \(\mu_S^{64}\), Verify가 만드는 32바이트
배열을 \(\mu_V^{32}\)라 하면
\[
  \operatorname{MuPrefix}_{32}(\mu_S^{64},\mu_V^{32})
  \;\Longleftrightarrow\;
  \forall i\in[0,32),\;\mu_S^{64}[i]=\mu_V^{32}[i].
\]
즉 \(\take_{32}(\mu_S^{64})=\mu_V^{32}\)이다.
\end{definition}

\begin{definition}[원시 분기(raw branch)와 유효 구간]
\label{def:raw-region}
64비트 word \(p,\ell\)에 대해 코드의 두 술어는
\[
  \operatorname{RawPreLen}(\ell)
    \Longleftrightarrow 0\le\operatorname{uint}(\ell)<2^{63},
  \qquad
  \operatorname{ValidRegion}(p,\ell)
    \Longleftrightarrow
    \operatorname{uint}(p)+\operatorname{uint}(\ell)\le 2^{64}
\]
이다. 첫 식은 Verify가 \identifier{prelen >> 63}으로 고르는 실제 분기에서
raw branch를 선택하게 하고, 둘째 식은 주소 덧셈의 wraparound를 배제한다.
\end{definition}

\subsection{Hoare 명제와 관계적 동치}

EasyCrypt의
\[
  \operatorname{hoare}[P:\Phi\Longrightarrow\Psi]
\]
는 \(P\)가 종료한다는 명제가 아니다. \(\Phi\)에서 시작하여 \textbf{종료한}
모든 실행이 \(\Psi\)를 만족한다는 부분정확성 명제다. 따라서 KeyGen의 retry
loop가 lossless인지, 거의 확실히 종료하는지, 논문과 같은 분포를 내는지는 별도
정리가 필요하다.

관계적 명제
\[
  \operatorname{equiv}[P\sim Q:\Phi\Longrightarrow\Psi]
\]
는 두 실행 상태의 사전관계 \(\Phi\)에서 시작한 \(P,Q\)의 결과가 사후관계
\(\Psi\)를 만족함을 뜻한다. 본 작업의 hash 정리는 SHAKE를 추상 함수로
치환하지 않고, 양쪽의 생성된 absorb/finalize/Keccak/squeeze 절차를 이
관계 논리로 직접 비교한다.

\subsection{정확 추적(exact trace)과 승인 꼬리(accepted tail)}

새로운 raw API 정리의 trace module은 실제 절차의 제어흐름을 복사하면서
\identifier{observed_mu}, pointer, length, \identifier{tail_reached} 같은 ghost
관찰값만 추가한다. exact equivalence는 실제 절차와 trace가 최종
\identifier{Glob.mem}과 반환값을 같게 만든다는 뜻이지, 모든 실행이 hash
tail에 도달한다는 뜻은 아니다.

Verify는 서명 길이, unpack, norm 등의 검사를 통과하기 전에 거절할 수 있다.
현재 정리는 이 분기를 숨기지 않고
\[
  \operatorname{VerifyAccept}\Longrightarrow\operatorname{TailReached}
\]
를 증명한 뒤, 성공한 실행에서만 실제 \identifier{__verify_hash_mu}에 전달된
VK/pre/message 주소와 길이를 관찰값에 연결한다. 이는 ``임의의 서명이 tail에
도달한다''는 주장과 구별해야 한다.

exact trace가 실제 절차와 같은 결과를 보존한다고 해서 trace 내부 hash 호출의
반환값과 종료 후 ghost 필드가 자동으로 관계 정리에 들어오는 것은 아니다.
현재 남은 \claim{OBL-MU-PRODUCT-REPLAY}는 두 completed sequential trace의
\identifier{observed_mu}를 기존 pRHL hash-return 정리에 연결하는 규칙이다.
EasyCrypt의 단순한 \identifier{sim}, \identifier{seq}, \identifier{byequiv}만으로
서로 다른 residual continuation을 제거할 수 없어 이 간선은
\statusdeferred 로 명시되었다.

\subsection{배열 크기와 논리적 길이}

추출된 이론은 여러 파라미터 집합을 담을 수 있는 고정 배열
\identifier{BArray2080}, \identifier{BArray2752},
\identifier{BArray2948}을 사용한다. 이것은 mode-2의 논리적 길이가 각각
2080, 2752, 2948이라는 뜻이 아니다. 실제 mode-2 API 길이는 검증키 992,
비밀키 1408, 서명 1474바이트이고, 현재 정리는 그중 첫 992바이트의 관계를
추적한다. 서명 codec 쪽에서는 별도로 challenge 32바이트와 low 1024바이트,
즉 서명의 첫 1056바이트 배치와 round trip을 추적한다. HBZ leaf 쪽의
\identifier{BArray8192}는 2048개의 32비트 계수를, symbol
\identifier{BArray2048}은 2048바이트를 담지만 mode-2 논리적 count는 1024다.
따라서 Week~8 frame 정리는 두 배열의 \([1024,2048)\) tail 보존을 별도로
명시한다.
